% intro.tex -- "Before you start": what is in the booklet, data, formulas, answer grid.
\unnumberedchapter[\faInfoCircle]{Before you start}

This booklet has 30 multiple-choice questions on \textbf{Module 5, Topic 3: Motion in
gravitational fields}. Each question is worth 1 mark, like Section I of the HSC exam. The
questions are new, but they test the same ideas as past HSC questions on this topic.

\begin{table}
  \centering
  \begin{tabular}{@{}L{22mm}L{82mm}L{30mm}@{}}
    \toprule
    \colhead{Section} & \colhead{Topic} & \colhead{Questions}\\
    \midrule
    Section A & Gravitational force and field strength & \hyperlink{pq.1}{1 to 6}\\
    Section B & Orbits and Kepler's laws & \hyperlink{pq.7}{7 to 18}\\
    Section C & Energy in orbits and escape velocity & \hyperlink{pq.19}{19 to 30}\\
    \bottomrule
  \end{tabular}
\end{table}

\textbf{How to use it}
\begin{itemize}
  \item Choose A, B, C or D. Write your answers in the grid below.
  \item Allow about 55 minutes for all 30 questions (about 1.8 minutes for each mark, as in
        the HSC exam). You can also do one section at a time.
  \item Use the data and formulas below. Show your working for the calculation questions.
  \item The answers are at the back. Each answer explains why the correct option is right
        \emph{and} why each wrong option is wrong. Click a question number to jump between a
        question and its answer.
\end{itemize}

\section*{Data}
$G = \qty{6.67e-11}{N.m^2.kg^{-2}}$ \qquad
$M_\mathrm{E} = \qty{6.0e24}{kg}$ \qquad
$r_\mathrm{E} = \qty{6.371e6}{m}$ \qquad
$g = \qty{9.8}{m.s^{-2}}$ (at Earth's surface)

\section*{Formulas used in this booklet}
\begin{formulatable}
  \frow{$F = \dfrac{GMm}{r^2}$}{gravitational force between two masses; $r$ is measured from
    centre to centre}
  \frow{$g = \dfrac{GM}{r^2}$}{gravitational field strength at a distance $r$ from the centre of
    a planet}
  \frow{$v = \dfrac{2\pi r}{T}$}{orbital speed from the orbital radius and the period}
  \frow[notonsheet]{$v = \sqrt{\dfrac{GM}{r}}$}{orbital speed in a circular orbit (from
    $\frac{mv^2}{r} = \frac{GMm}{r^2}$)}
  \frow{$\dfrac{r^3}{T^2} = \dfrac{GM}{4\pi^2}$}{Kepler's third law, for any object in orbit
    around a mass $M$}
  \frow{$U = -\dfrac{GMm}{r}$}{gravitational potential energy; zero at infinity, negative
    everywhere else}
  \frow[notonsheet]{$E = -\dfrac{GMm}{2r}$}{total energy (kinetic + potential) of a satellite in a
    circular orbit; it is half of $U$}
  \frow{$v_\mathrm{esc} = \sqrt{\dfrac{2GM}{r}}$}{escape velocity at a distance $r$ from the centre
    of a planet}
\end{formulatable}
Remember: $r$ is the distance from the \emph{centre} of the planet. For a satellite,
$r = r_\mathrm{E} + \text{altitude}$.

\section*{Answer grid}
\begin{table}
  \centering\renewcommand\arraystretch{1.6}
  \begin{tabular}{@{}l*{10}{C{9mm}}@{}}
    \toprule
    \colhead{Question} & 1 & 2 & 3 & 4 & 5 & 6 & 7 & 8 & 9 & 10\\
    Answer & & & & & & & & & & \\
    \midrule
    \colhead{Question} & 11 & 12 & 13 & 14 & 15 & 16 & 17 & 18 & 19 & 20\\
    Answer & & & & & & & & & & \\
    \midrule
    \colhead{Question} & 21 & 22 & 23 & 24 & 25 & 26 & 27 & 28 & 29 & 30\\
    Answer & & & & & & & & & & \\
    \bottomrule
  \end{tabular}
\end{table}
\begin{center}
  Score: \underline{\hspace{14mm}} / 30
\end{center}
